\section{A common-base ladder determinant with an unresolved odd-prime cost}
\label{sec:common-base}

Approximating the terms of a functional identity separately can destroy
its arithmetic advantage: eliminating their unwanted constants multiplies
large coefficients. Here the polylogarithm arguments remain powers of
one variable throughout the construction. This gives a determinant whose
analytic decay beats its entire cost at the prime 2. The missing step is
the other primes, not nonvanishing. We develop record \recordid{L03}
and then a related signed-integral alternative.

\subsection{A rational-coefficient germ and its endpoint determinant}
Put $q=(1-i)/2$, $z_m=-q^m$ for $m=2,3,5$, and
$(c_2,c_3,c_5)=(225/136,-25/34,9/68)$. Write
\[
 \mathcal I_s(z_mt^m)=
  \frac{\Li_s(z_mt^m)-\Li_s(\overline z_mt^m)}{2i},
 \qquad L(t)=-\log(1-t^2/2).
\]
This defines an analytic germ, not the imaginary part operation on
complex $t$. All its Taylor coefficients are rational. Define
\begin{equation}\label{eq:common-H}
 H(t)=\sum_{m=2,3,5}c_m\left(
 \mathcal I_4(z_mt^m)+\frac m2L(t)\mathcal I_3(z_mt^m)
              +\frac{m^2}{12}L(t)^2\mathcal I_2(z_mt^m)\right)
       =\sum_{j\geq1}h_jt^j.
\end{equation}
The completed Broadhurst identity~\eqref{eq:ladder} gives $H(1)=\theta$.
Its use here has the same external dependency as Section~\ref{sec:mellin}.
Each $h_j$ is obtained by finite convolution of the displayed series,
so the following polynomial is reproducible over $\Q$:
\begin{equation}\label{eq:common-J}
 b_i=h_i,\quad T_{ij}=h_{i+j}-h_{i+j-1}\quad(1\leq i,j\leq n),
 \qquad
 J_n(X)=\det\begin{pmatrix}X&-b^T\\b&T\end{pmatrix}.
\end{equation}
For $n=0$ put $J_0(X)=X$. Its degree is at most one, even though its
matrix size grows. With
$c_j(X)=X-\sum_{\ell=1}^jh_\ell$, adjacent row and column differences
give the exact identity
\begin{equation}\label{eq:common-hankel}
 \det[c_{i+j}(X)]_{0\leq i,j\leq n}=(-1)^nJ_n(X),
 \qquad \sum_{j\geq0}c_j(X)t^j=\frac{X-H(t)}{1-t}.
\end{equation}
No invertibility of $T$ or Pad\'e normality is required.

\subsection{Nonvanishing from a logarithmic branch}
As $t$ increases to $\sqrt2$, the polylogarithm arguments tend to
$i,e^{i\pi/4},e^{-i\pi/4}$, none equal to 1. The coefficient of
$L(t)^2$ at that endpoint is
\[
 A_+=\frac{75}{272}\bigl(2\beta(2)-3\operatorname{Cl}_2(\pi/4)\bigr),
 \qquad \operatorname{Cl}_2(\phi)=\sum_{j\geq1}\frac{\sin(j\phi)}{j^2}.
\]
The alternating series gives $\beta(2)<1$. The first three Clausen
terms and the absolute tail bound $\sum_{j\geq4}j^{-2}<1/3$ give
\[
 \operatorname{Cl}_2(\pi/4)>
       \frac{5\sqrt2}{9}+\frac14-\frac13>\frac{25}{36},
 \qquad A_+<-\frac{25}{1088}.
\]
Thus $H$ has a genuine logarithm-squared singularity. Subtracting any
constant $X$ and dividing by $1-t$ cannot remove it. Kronecker's
classical Hankel criterion says that a power series is rational exactly
when all its sufficiently large leading Hankel determinants vanish
\cite[Theorem 1]{KroneckerNote}. Applied to~\eqref{eq:common-hankel},
it proves that $J_n(X)\ne0$ for infinitely many $n$ for every fixed
complex $X$, in particular $X=\theta$. No irrationality assumption or
positivity argument enters this step.

\subsection{A common analytic domain and a determinant bound}
Let $E=t\,d/dt$. The rational function $U$ in~\eqref{eq:U} satisfies
\[
 U(t)=\sum_m c_m m^4\mathcal I_0(z_mt^m).
\]
Define $Y_0=U$, $EY_j=Y_{j-1}$ and $Y_j(0)=0$ for $1\leq j\leq4$.
Then
\begin{equation}\label{eq:common-tower}
 H=Y_4+\tfrac12LY_3+\tfrac1{12}L^2Y_2,\qquad
 Y_j(t)=\frac1{(j-1)!}\int_0^1U(tx)(-\log x)^{j-1}\frac{\dd x}{x}.
\end{equation}
These identities follow coefficientwise. They also specify analytic
continuation by integration on radial segments.

For $p_2=t^4+4$, $p_3=t^6-4t^3+8$, $p_5=t^{10}-8t^5+32$,
let $D=p_2p_3p_5/(t^2+2t+2)^2$. Polynomial division shows
\[
 (2-t^2)D(t)\mid t^{120}-2^{60}.
\]
All possible finite singularities lie at radius $\sqrt2$ on the
120th-root rays. Remove those rays outward from that radius. The
remaining domain is star-shaped, and~\eqref{eq:common-tower} gives
single-valued holomorphic functions there. There is no singularity at
zero because $U=O(t^2)$.

In the reciprocal variable put
\[
 G(w)=\frac{\theta-H(1/w)}{w-1},\qquad
 K=\{w:w^{120}\in[0,2^{-60}]\}.
\]
The apparent pole at $w=1$ is removable precisely because $H(1)=\theta$.
Thus $G$ is holomorphic off the compact star $K$ and vanishes at infinity.
The logarithmic capacity of a segment of length $a$ is $a/4$; the
monic polynomial-preimage rule gives
\[
 \operatorname{cap}(K)=\bigl(2^{-60}/4\bigr)^{1/120}=2^{-31/60}.
\]
For this particular preimage the rule also follows by composing the
segment's Green function with $w^{120}$ and dividing by 120.

P\'olya's determinant inequality~\cite[Theorem 2.1, one variable]{Polya}
now yields
\begin{equation}\label{eq:common-real}
 \log|J_n(\theta)|\leq-\frac{31}{60}(\log2)n^2+o(n^2),
\end{equation}
with $\log0=-\infty$. To check the exponent normalization, the Laurent
coefficients of $G$ are $c_j(\theta)$. Their Hankel determinant is
\[
 \frac1{(n+1)!(2\pi i)^{n+1}}\int_{\Gamma^{n+1}}
       \prod_jG(w_j)\prod_{i<j}(w_j-w_i)^2\prod_j\dd w_j.
\]
Here $\Gamma$ surrounds $K$ in its analytic complement.
The maximum squared Vandermonde on successively smaller contour
neighborhoods has leading logarithm
$n(n+1)\log\operatorname{cap}(K)$. The remaining factors contribute
only $O(n\log n)$ for a fixed contour. This explains why there is one
capacity exponent on the $n^2$ scale, not twice that exponent.

\subsection{The complete leading cost at 2}
For $J_n\ne0$ write $\Prim(J_n)=\kappa_nJ_n$.
We give a short bound sufficient for the leading rate; the stronger
finite-index residue-rank correction is unnecessary here.
In $\Q_2(i)$ take $\pi_2=1-i$, with $v_{\pi_2}(2)=2$.
Substitute $x=qt$, so $v_{\pi_2}(q)=-1$. The polylogarithm blocks become
\[
 \mathcal B_{m,s}(x)=\frac{\Li_s(-x^m)-\Li_s(-i^mx^m)}{2i},\qquad
 [x^{mj}]\mathcal B_{m,s}
       =\frac{(-1)^j(1-i^{mj})}{2ij^s}.
\]
For $m=2$ these vanish at even $j$. For $m=3,5$ they vanish when
$4\mid j$; at odd $j$ the valuation is $-1$, and at $j\equiv2\pmod4$
it is $-2s$. Including $c_m,m/2,m^2/12$, multiplication by $-64$
makes each of the three blocks in~\eqref{eq:common-H} integral.
The worst valuations before that multiplication, for each of the
three blocks, are $-6,-10,-12$ at $m=2,3,5$.

In the $x$ variable $L=-\log(1-ix^2)$. Through degree $2n$ its
coefficients have 2-adic valuation at least $-r$, where
$r=\lfloor\log_2 n\rfloor$; those of $L^2$ have valuation at least
$-2r$. Hence $f_j=\lambda q^{-j}h_j$ is integral for $j\leq2n$,
where $\lambda=-64\,2^{2r}$. Scale row and column $i$ of the bordered
matrix by $q^{-i}$, with index zero for the border, and scale the whole
matrix by $\lambda$. Its entries are $\lambda X$, $\pm f_i$ and
$f_{i+j}-q^{-1}f_{i+j-1}$, all integral. The determinant multiplier has
2-adic valuation
\[
 B_n=\frac{n(n+1)}2+(6+2r)(n+1).
\]
Since the original coefficients are rational, their valuations over
$\Q_2(i)$ are twice their ordinary 2-adic valuations. We conclude
\begin{equation}\label{eq:common-two}
 v_2(\kappa_n)\leq B_n=\tfrac12n^2+O(n\log n).
\end{equation}
Combining~\eqref{eq:common-real} and~\eqref{eq:common-two} gives
\[
 \log|2^{B_n}J_n(\theta)|
       \leq-\frac{\log2}{60}n^2+o(n^2).
\]
These are shrinking, infinitely often nonzero linear forms with
2-integral coefficients. They are not necessarily integer forms.

\subsection{The precise missing inequality}
Put $d_n=2^{\max(0,v_2(\kappa_n))}$. It would suffice to prove
\begin{equation}\label{eq:common-missing}
 \limsup_{\substack{n\to\infty\\J_n\ne0}}
       \frac{\log(\kappa_n/d_n)}{n^2}<\frac{\log2}{60}.
\end{equation}
This includes all odd-prime costs and any compensating numerator
content; it is a sufficient target, not a necessary condition for every
use of the germ. No bound of this strength is established. Finite
primitive failures at $n=12,16,20$ do not settle the asymptotic question.
To reproduce a form, expand~\eqref{eq:common-H} through degree $2n$,
build~\eqref{eq:common-J} over $\Q[X]$, and clear its two coefficients
followed by their full gcd. This procedure does not require solving a
possibly singular Pad\'e system.

\subsection{A signed-integral alternative with only one LCM layer}
The same rational kernel also gives a different, completely explicit
starting point. Retain $D$ above and put $N_0=17DU\in\Z[x]$; it has
a zero of order two at zero. For arbitrary $S\in\Z[x]$ and $A\in\Z$ set
\[
 F=x(1-x)^2N_0^4S,\qquad
 V=17\frac{E^3F}{N_0},\qquad P=A+DV.
\]
At most three differentiations can remove at most three of the four
$N_0$ factors, so $V\in\Z[x]$. With $N=DU=N_0/17$ one has
$NV=E^3F$. If $F=\sum f_jx^j$, the constraints of
\eqref{eq:constraints} become $F(1)=EF(1)=0$, which hold by construction.
Consequently~\eqref{eq:mellin-pair} gives exactly
\begin{equation}\label{eq:signed-rodrigues}
 I(P)=A\theta+\sum_j\frac{f_j}{j}.
\end{equation}
Thus $\operatorname{lcm}(1,\ldots,\deg F)$ clears the rational remainder,
rather than its fourth power. This is an arithmetic improvement, not
a smallness or nonvanishing theorem.

There are infinitely many nonzero source polynomials giving the zero
form. For any nonzero $Z\in\Z[x]$, let
\[
 R_0=N_0^5(1-x)^3Z,\qquad F=ER_0,
 \qquad S=(5N_0'(1-x)-3N_0)Z+N_0(1-x)Z'.
\]
Then this $F$ has the prescribed form and
$\sum f_j/j=R_0(1)-R_0(0)=0$. With $A=0$ the source is nonzero,
since $E^4$ is injective on polynomials with zero constant coefficient.
A small-vector argument must therefore distinguish the kernel of the
map to the final two coefficients from genuine small nonzero forms.

There is also a limitation on an absolute-integral estimate. In the
notation of Section~\ref{sec:mellin}, the constraint preserves
$\int PU(-\log x)\dd x/x=A M_1$, where $M_1>9/10$ is proved there.
As $W(\rho)\geq(\log2)^2\rho/12$,
\[
 \int_0^1|PU|W(-\log x)\frac{\dd x}{x}
       \geq\frac{(\log2)^2}{12}|A|M_1.
\]
After a positive normalization making the nonzero target coefficient
an integer this is greater than $1/30$.
Hence a successful signed choice must exploit cancellation in the
actual integral and establish nonvanishing separately. Formula
\eqref{eq:signed-rodrigues} supplies free parameters and exact arithmetic;
it does not yet supply that choice or those two proofs.
